% TOP_PROBLEM: {"id": 1900020, "title": "Geometry of Continued Fractions — Further open questions", "category_id": 6, "category": {"id": 6, "name": "geometry", "display_name": "Geometry", "description": "Euclidean and non-Euclidean geometry, geometric structures.", "slug": "geometry", "order_index": 6, "created_at": "2026-07-31T15:26:25.671Z"}, "status": "partially_solved"}
% FIRST_POSTED: null
% ENTRY_KIND: statement_only
% Imported from: batch07.tex; UnsolvedMath ID 1900020
\documentclass[11pt]{book}
\usepackage{fontspec}
\setmainfont{Latin Modern Roman}
\setsansfont{Latin Modern Sans}
\usepackage{amsmath,amssymb,mathtools}
\usepackage{unicode-math}
\setmathfont{Latin Modern Math}
\usepackage{xurl}
\usepackage[hidelinks]{hyperref}
\newcommand{\sref}[2]{\hyperlink{source:#1:#2}{[S#2]}}
\newcommand{\cataloguescope}[1]{\paragraph{Mathematical question.}#1}
\begin{document}
% UnsolvedMath ID 1900020; source catalogue code AMR-018-0020
\setcounter{section}{1900019}
\section{Geometry of Continued Fractions — Further open questions}\label{1900020}
{\small\sffamily UnsolvedMath ID 1900020 · Geometry}
\subsection{Definitions and mathematical statement}

\paragraph{Shared notation.}$\mathbb N=\{1,2,\ldots\}$, and $\mathbb Z,\mathbb Q,\mathbb R,\mathbb C$ denote the integers, rationals, reals and complex numbers. Any other conventions are specified locally.


Let $K$ be a degree-three extension of $\mathbb Q$ all of whose embeddings into $\mathbb C$ have image in $\mathbb R$, and fix one real embedding to compare its elements with $0$ and $1$. Let $y,z\in K$ satisfy $0<y,z<1$ and suppose that $1,y,z$ are linearly independent over $\mathbb Q$. Define the classical Jacobi--Perron iteration by $v_0=(x_0,y_0,z_0)=(1,y,z)$ and
\[
 a_i=\left\lfloor\frac{z_i}{y_i}\right\rfloor,\qquad
 b_i=\left\lfloor\frac{x_i}{y_i}\right\rfloor,\qquad
 (x_{i+1},y_{i+1},z_{i+1})=(y_i,z_i-a_i y_i,x_i-b_i y_i).
\]
Here $\lfloor u\rfloor$ denotes the greatest integer at most $u$. Each step is an invertible integer linear transformation, so linear independence is preserved and none of the coordinates vanishes. The output is the sequence of integer pairs $(a_i,b_i)$. It is \emph{eventually periodic} if there are integers $N\ge0$ and $p\ge1$ such that $(a_{i+p},b_{i+p})=(a_i,b_i)$ for all $i\ge N$. \sref{1900020}{3}
\paragraph{Statement recorded in the supplied source.}
For every such field $K$ and every such pair $y,z$, is this Jacobi--Perron output eventually periodic? \sref{1900020}{2}\sref{1900020}{3}
\subsection{Short English statement}
Does the classical Jacobi--Perron algorithm eventually repeat for every independent basis triple of a totally real cubic field with the two nonconstant entries between zero and one?
\subsection{Sources}
\begin{description}
\item[\hypertarget{source:1900020:1}{S1}] ULAM.AI and the UnsolvedMath Contributors, \emph{UnsolvedMath: A Curated Collection of Open Mathematics Problems}, record 1900020 (AMR-018-0020). CC BY 4.0. \url{https://huggingface.co/datasets/ulamai/UnsolvedMath}.
\item[\hypertarget{source:1900020:2}{S2}] Oleg Karpenkov, \emph{Open problems in geometry of continued fractions}, 2017. Problem 20 and its surrounding definitions and remarks. \url{https://arxiv.org/pdf/1712.01450}.
\item[\hypertarget{source:1900020:3}{S3}] Oleg Karpenkov, \emph{On Hermite's problem, Jacobi--Perron type algorithms, and Dirichlet groups}, arXiv:2101.12707, 2021. Section 2.1, pp. 5--7, algorithm and Problem 1 (Jacobi's Last Theorem). \url{https://arxiv.org/pdf/2101.12707}.
\end{description}
\label{end:1900020}
\end{document}
